\documentclass[11pt,a4paper]{article}

\usepackage[utf8]{inputenc}
\usepackage[margin=0.9in]{geometry}
\usepackage{amsmath,amssymb,amsthm,amsfonts}
\usepackage{mathtools}
\usepackage{hyperref}
\usepackage{booktabs}
\usepackage{cite}
\usepackage{microtype}
\usepackage{array}

\hypersetup{
    colorlinks=true,
    linkcolor=blue,
    citecolor=red,
    urlcolor=blue
}

% Theorem Environments
\theoremstyle{plain}
\newtheorem{theorem}{Theorem}[section]
\newtheorem{lemma}[theorem]{Lemma}
\newtheorem{proposition}[theorem]{Proposition}
\newtheorem{corollary}[theorem]{Corollary}
\newtheorem{conjecture}[theorem]{Conjecture}
\newtheorem{problem}[theorem]{Problem}

\theoremstyle{definition}
\newtheorem{definition}[theorem]{Definition}
\newtheorem{assumption}[theorem]{Assumption}
\newtheorem{example}[theorem]{Example}

\theoremstyle{remark}
\newtheorem{remark}[theorem]{Remark}

\numberwithin{equation}{section}

\title{\textbf{Continuous Pascal Simplexes, Sierpi\'nski Gasket Laplacians, and Multifractal Singularity Spectra}}
\author{\textbf{Reinaldo Maia Silva-Filho} \\ \textit{Universidade Federal de Lavras (UFLA)}}
\date{\today}


\DeclareMathOperator*{\argmin}{argmin}
\providecommand{\Hyp}{\mathbb{H}}
\providecommand{\Sph}{\mathbb{S}}
\providecommand{\II}{\mathrm{II}}
\providecommand{\R}{\mathbb{R}}
\providecommand{\Area}{\mathrm{Area}}
\providecommand{\Hsn}{\mathcal{H}}
\providecommand{\BV}{\mathrm{BV}}
\providecommand{\TV}{\mathrm{TV}}
\providecommand{\cutnorm}[1]{\|#1\|_{\square}}

\begin{document}

\maketitle

\begin{abstract}
We examine connections between the analytic continuation of Pascal's $m$-simplex and analysis on post-critically finite (p.c.f.) self-similar fractals. We recall Kigami's Laplacian on the Sierpi\'nski $m$-simplex, derive its renormalization factor $m+3$ from the level-one network, and recall the resulting spectral dimension $d_s = \frac{2\ln(m+1)}{\ln(m+3)}$ (Kigami--Lapidus). We then pose, as an explicitly specified open problem, whether non-local forms built from the Beta-kernel of Chapter~3 converge to Kigami's form; the naive construction is ill-defined, for reasons we make explicit. Using the Barnes $G$-function representation of the continuous row entropy, we prove its asymptotic expansion $\mathcal{E}(x) = \frac{x^2}{2} - \frac{x}{2}\ln x + (1 - \frac{1}{2}\ln 2\pi)x - \frac16\ln x + 2\zeta'(-1) - \frac1{12} - \frac{1}{180x^2} + \mathcal{O}(x^{-4})$, and we discuss a heuristic parabolic multifractal parametrization. Finally, we recall the classical fact that Pascal's triangle modulo $2$ has the Sierpi\'nski gasket as its scaling limit, with box-counting dimension $\ln 3/\ln 2$.
\end{abstract}

\tableofcontents
\newpage

% ==============================================================================
\section{Introduction and Historical Context}
% ==============================================================================

The emergence of fractal structures from combinatorics is historically rooted in Lucas' Theorem (1878) \cite{lucas1878}, which states that for prime $p$, the binomial coefficient $\binom{n}{k} \pmod p$ is non-zero if and only if the base-$p$ digits of $k$ do not exceed those of $n$. In the simplest case $p=2$, the distribution of odd binomial coefficients in Pascal's triangle converges in the Hausdorff metric to the \textbf{Sierpi\'nski Gasket} $K \subset \mathbb{R}^2$ with Hausdorff dimension:
\begin{equation}
d_H = \frac{\ln 3}{\ln 2} \approx 1.58496.
\end{equation}
More generally, the Pascal $m$-pyramid modulo $2$ (layer index together with $m-1$ free multinomial indices) generates the \textbf{Sierpi\'nski $m$-Simplex} in $\mathbb{R}^{m}$, with $m+1$ self-similar pieces and Hausdorff dimension $d_H = \frac{\ln(m+1)}{\ln 2}$.

Traditional analysis on fractals, pioneered by Jun Kigami and Robert Strichartz \cite{kigami2001, strichartz2006}, constructs Laplacians on self-similar sets via discrete graph Laplacians on sequence networks:
\begin{equation}
\Delta_{\text{Kigami}} u = \lim_{k \to \infty} r_m^k \Delta_k u,
\end{equation}
where $r_m = m+3$ is the harmonic renormalization factor: the cell count $m+1$ times the resistance scaling $\rho = \frac{m+3}{m+1}$ (Proposition~\ref{assump:kigami_decimation}; for $m=2$, $r_2 = 5$).

In this work, expanding the continuous simplicial framework established in our foundational series \cite{silvafilho2026simplex, silvafilho2026waves, silvafilho2026inter}, we examine which links between the continuation of the Pascal simplex via Euler's Gamma function $\Gamma(z)$ and analysis on fractals can be proved. Our findings are mixed. The spectral dimension of the Sierpi\'nski simplex is Kigami's classical result, and we only recall it. A derivation of Kigami's Laplacian from the Beta-kernel operators of Chapter~3 remains open (Problem~\ref{thm:kigami_convergence}). The new result proved here is the asymptotic expansion of the continuous row entropy (Proposition~\ref{thm:multifractal}). The multifractal parametrization is a heuristic.

% ==============================================================================
\section{Beta-Kernel Decimations and Kigami's Fractal Laplacian}
% ==============================================================================

Let $K \subset \mathbb{R}^{m}$ be the regular Sierpi\'nski $m$-simplex generated by the Iterated Function System (IFS) $\mathcal{F} = \{F_1, \dots, F_{m+1}\}$ with contraction factor $r = 1/2$:
\begin{equation}
F_j(\mathbf{x}) = \frac{1}{2}\mathbf{x} + \frac{1}{2}\mathbf{v}_j, \quad j=1,\dots,m+1,
\end{equation}
where $\{\mathbf{v}_j\}_{j=1}^{m+1}$ are the vertices of a regular $m$-simplex.

\begin{definition}[Decimated Simplicial Fractional Operator]
For refinement level $k \in \mathbb{N}$ and fractional order $\alpha > 0$, we consider the \emph{formal} decimated operator on $L^2(K, \mu)$ (where $\mu$ is the normalized Hausdorff measure)
\begin{equation}
\label{eq:decimated_op}
\mathcal{L}_k^\alpha u(\mathbf{x}) \coloneqq (m+3)^k \sum_{|w|=k} \Delta_{\Delta_m}^{\alpha \cdot 2^{-k}} \left(u \circ F_w\right)\left(F_w^{-1}(\mathbf{x})\right) \cdot \chi_{F_w(K)}(\mathbf{x}),
\end{equation}
where $w = (w_1, \dots, w_k) \in \{1, \dots, m+1\}^k$ denotes word sequences, and $\Delta_{\Delta_m}^\alpha$ is the Fractional Simplicial Laplacian. Chapter~3 defines $\Delta^\alpha_{\Delta_m}$ on $L^2(\mathbb{R}^{m-1})$, with kernel on $\Delta_{m-1}(\alpha) \subset \mathbb{R}^{m-1}$. Since $K \subset \mathbb{R}^m$, here we use the same construction one dimension higher, with the kernel on $\Delta_m(\alpha) \subset \mathbb{R}^m$ ($m+1$ parts), acting on functions on $\mathbb{R}^m$. The expression is formal, because such an operator acts on functions on $\mathbb{R}^m$, not on $K$ (see Remark~\ref{rem:naive_fails}).
\end{definition}

\begin{proposition}[Harmonic Decimation Renormalization Factor]
\label{assump:kigami_decimation}
Let $\Gamma_0$ be the complete graph on the vertices $p_1, \dots, p_{m+1}$ of $K$ with unit conductances. Let $\Gamma_1$ be the level-one network: its vertices are the $p_i$ and the midpoints $p_{ij} = p_{ji}$ ($i \ne j$), and each cell $F_i(K)$ carries a complete graph with unit conductances on $\{p_i\} \cup \{p_{ij} : j \ne i\}$. The trace (Schur complement) of $\Gamma_1$ on $\{p_1, \dots, p_{m+1}\}$ is the complete graph with conductances $\frac{m+1}{m+3}$. Hence the resistance scaling is $\rho = \frac{m+3}{m+1}$, and the Laplacian renormalization factor with respect to the self-similar measure (weights $\frac{1}{m+1}$) is $r_m = (m+1)\rho = m+3$. For $m = 2$, $r_2 = 5$.
\end{proposition}
\begin{proof}
By symmetry the trace is a complete graph with a common conductance $c$. Impose $u(p_1) = 1$ and $u(p_j) = 0$ for $j \ne 1$, and extend $u$ harmonically. By symmetry $u(p_{1j}) = a$ for $j \ne 1$ and $u(p_{jk}) = b$ for $j, k \ne 1$. The vertex $p_{1j}$ has $m$ neighbours in each of the cells $1$ and $j$, so harmonicity gives $2ma = 1 + (m-1)a + 0 + (m-1)b$. The vertex $p_{jk}$ gives $2mb = 2a + 2(m-2)b$. Hence $b = a/2$ and $a = \frac{2}{m+3}$. The current leaving $p_1$ is $m(1 - a) = \frac{m(m+1)}{m+3}$, which equals $mc$ for the trace network, so $c = \frac{m+1}{m+3}$. This is the classical computation behind Kigami's harmonic structure on the Sierpi\'nski simplex \cite{kigami2001}.
\end{proof}

\begin{problem}[Beta-Kernel Forms and Kigami's Laplacian]
\label{thm:kigami_convergence}
Let $\widetilde{\mathcal{K}}_\alpha$ be the normalized Beta-kernel on $\Delta_m(\alpha) \subset \mathbb{R}^m$, let $\kappa_\alpha(\mathbf{z}) \coloneqq \frac12\big(\widetilde{\mathcal{K}}_\alpha(\mathbf{z}) + \widetilde{\mathcal{K}}_\alpha(-\mathbf{z})\big)$ be its symmetrization, and for $u \in L^2(K, \mu)$ define
\begin{equation}
\label{eq:beta_forms}
\mathcal{E}_k(u) \coloneqq c_k \iint_{K \times K} \kappa_{\alpha_k}\big(2^k(\mathbf{x} - \mathbf{y})\big)\, |u(\mathbf{x}) - u(\mathbf{y})|^2 \, d\mu(\mathbf{x})\, d\mu(\mathbf{y}).
\end{equation}
Do there exist orders $\alpha_k > 0$ and constants $c_k > 0$ such that $\mathcal{E}_k$ converges in the sense of Mosco on $L^2(K, \mu)$ to Kigami's energy form $\mathcal{E}_K$, so that the associated generators converge to $\Delta_{\text{Kigami}}$ in the strong resolvent sense?
\end{problem}

The forms \eqref{eq:beta_forms} are integrated against $\mu$ and evaluate $u$ only on $K$, so they avoid the defects listed below. Since $\kappa_{\alpha_k}$ is bounded, each $\mathcal{E}_k$ is a bounded quadratic form on $L^2(K,\mu)$. The kernel $\kappa_{\alpha_k}(2^k(\mathbf{x}-\mathbf{y}))$ vanishes unless $2^k(\mathbf{x}-\mathbf{y}) \in \Delta_m(\alpha_k) \cup (-\Delta_m(\alpha_k))$, so only pairs at distance of order $\alpha_k 2^{-k}$ interact, and these include pairs in neighbouring level-$k$ cells. The coupling across cells is essential. If the integral were restricted to pairs lying in the same level-$k$ cell $F_w(K)$, then for $u = \mathbf{1}_{F_1(K)}$ every $\mathcal{E}_k(u)$ would vanish, since each level-$k$ cell with $k \ge 1$ lies in a single level-one cell. But $u$ is not in the domain of $\mathcal{E}_K$: for the gasket its level-$n$ renormalized graph energies equal $4(5/3)^n$. The Mosco inequality $\liminf_k \mathcal{E}_k(u_k) \ge \mathcal{E}_K(u)$ would then fail for $u_k = u$ and every choice of $\alpha_k, c_k$. For the forms \eqref{eq:beta_forms}, by contrast, $\mathcal{E}_k(\mathbf{1}_{F_1(K)}) > 0$, because pairs on opposite sides of the junction points of $F_1(K)$ interact.

The Beta-kernel enters through the fixed jump kernel $\kappa_{\alpha_k}$, which may not be replaced by an arbitrary one. Without this requirement the question would be answered by Kigami's own discrete approximations. For fixed $\alpha_k = \alpha$, a scaling count suggests $c_k \asymp \big((m+1)(m+3)\big)^k$: the factor $(m+3)^k$ is the renormalization of Proposition~\ref{assump:kigami_decimation}, and $(m+1)^k$ compensates the $\mu \otimes \mu$ mass, of order $(m+1)^{-k}$, of the pairs at distance $\mathcal{O}(2^{-k})$. Non-local approximating forms with indicator kernels $\mathbf{1}_{\{|\mathbf{x}-\mathbf{y}|<r\}}$ and their subsequential $\Gamma$-limits on fractals are studied by Kumagai and Sturm \cite{kumagaisturm2005}. Two points separate Problem~\ref{thm:kigami_convergence} from that setting: Mosco convergence of the whole sequence to $\mathcal{E}_K$ itself is required, not comparability with $\mathcal{E}_K$ or convergence of a subsequence; and the support of $\kappa_\alpha$ is not invariant under the symmetry group of the simplex, so the edge directions of $K$ are weighted differently while $\mathcal{E}_K$ is invariant. We have no evidence either way.

\begin{remark}[Why the naive construction fails]
\label{rem:naive_fails}
The direct choice, with $\alpha_k \to d_H - 1$ and the energy forms
$\mathcal{E}_k(u,u) = (m+3)^k \sum_{|w|=k} c_k \int_{F_w(K)} d\mathbf{x}\int_{\Delta_{m}(\alpha 2^{-k})} \binom{\alpha 2^{-k}}{\mathbf{y}}|u(\mathbf{x}+\mathbf{y}) - u(\mathbf{x})|^2 d\mathbf{y}$, with Lebesgue measure $d\mathbf{x}$ on $\mathbb{R}^m$ and the kernel on $\Delta_m(\alpha 2^{-k}) \subset \mathbb{R}^m$. That construction does not work, for three reasons.
(a) $K$ has Lebesgue measure zero, so $\int_{F_w(K)} d\mathbf{x} = 0$ and the forms vanish identically.
(b) For $\mathbf{x} \in K$ the translate $\mathbf{x} + \mathbf{y}$ generally leaves $K$, where $u$ is undefined.
(c) Nothing forces the weights to converge to the harmonic factors, and the choice $\alpha_k \to d_H - 1$ has no justification.
Any positive answer requires a construction on $(K,\mu)$ itself, such as \eqref{eq:beta_forms}, for example via Mosco convergence as in \cite{kigami2001}. The renormalization factor $m+3$ is then given by Proposition~\ref{assump:kigami_decimation}.
\end{remark}

\begin{theorem}[Spectral Dimension of the Sierpi\'nski Simplex; Kigami--Lapidus \cite{kigamilapidus1993}, see also \cite{kigami2001}]
The eigenvalue counting function of Kigami's Laplacian on the Sierpi\'nski $m$-simplex satisfies the fractal Weyl law
\begin{equation}
c_1 \lambda^{d_s / 2} \le N(\lambda) \le c_2 \lambda^{d_s / 2} \quad \text{for } \lambda \ge 1,
\end{equation}
where the spectral dimension $d_s$ is given analytically by
\begin{equation}
d_s = \frac{2 d_H}{d_w} = \frac{2\ln(m+1)}{\ln(m+3)}, \qquad d_w = \log_2(m+3).
\end{equation}
with constants $0 < c_1 \le c_2$. In general $N(\lambda)/\lambda^{d_s/2}$ has no limit: for the gasket it oscillates log-periodically. This is a classical result about Kigami's Laplacian. It does not depend on Problem~\ref{thm:kigami_convergence}.
\end{theorem}

% ==============================================================================
\section[Multifractal Formalism via Barnes G-Function]{Multifractal Formalism \\ via the Barnes \texorpdfstring{$G$}{G}-Function}
% ==============================================================================

In the multifractal formalism of Halsey et al.\ \cite{halsey1986} (see also \cite{falconer2014}), the $L^q$-scaling exponents $\tau(q)$ and the singularity spectrum $f(\alpha)$ are defined for a \emph{measure}. One would like to define them through the moments $\int_0^x \binom{x}{y}^q dy$, but no measure with these moments has been constructed. We therefore prove only an asymptotic statement about the derivative of these moments at $q = 0$, which is the row entropy $\mathcal{E}(x)$, and keep the multifractal reading as a heuristic.

\begin{proposition}[Asymptotics of the Continuous Row Entropy]
\label{thm:multifractal}
The continuous row entropy $\mathcal{E}(x) \coloneqq \int_0^x \ln\binom{x}{y}\,dy = x(x+1) - x\ln(2\pi) - x\ln\Gamma(x+1) + 2\ln G(x+1)$ (Chapter~3) satisfies, as $x \to \infty$,
\begin{equation}
\mathcal{E}(x) = \frac{1}{2}x^2 - \frac{1}{2}x\ln x + \left(1 - \frac{1}{2}\ln(2\pi)\right)x - \frac{1}{6}\ln x + 2\zeta'(-1) - \frac{1}{12} - \frac{1}{180\,x^2} + \mathcal{O}(x^{-4}).
\end{equation}
In particular $\mathcal{E}(x)/x^2 \to \frac12$, which is the integral $\int_0^1 H(t)\,dt$ of the binary entropy in nats.
\end{proposition}

\begin{proof}
Insert Stirling's series $\ln\Gamma(x+1) = x\ln x - x + \frac{1}{2}\ln(2\pi x) + \frac{1}{12x} - \frac{1}{360x^3} + \mathcal{O}(x^{-5})$ and the Barnes asymptotic series $\ln G(x+1) = \frac{x^2}{2}\ln x - \frac{3x^2}{4} + \frac{x}{2}\ln(2\pi) - \frac{1}{12}\ln x + \zeta'(-1) - \frac{1}{240x^2} + \mathcal{O}(x^{-4})$ \cite{barnes1900}, whose $x^{-2}$ coefficient is $B_4/8$. The $x^2\ln x$ terms cancel between $-x\ln\Gamma(x+1)$ and $2\ln G(x+1)$. The term $\frac{1}{12x}$ contributes $-\frac{1}{12}$ at order $x^0$. At order $x^{-2}$, $-x \cdot (-\frac{1}{360x^3}) + 2 \cdot (-\frac{1}{240x^2}) = -\frac{1}{180x^2}$. Collecting the remaining terms gives the expansion. Numerically, $x^2$ times the difference between $\mathcal{E}(x)$ and the terms up to order $x^0$ tends to $-0.005556 = -1/180$. Replacing $-\frac16\ln x$ by $-\frac1{12}\ln x$ leaves an error of $0.77$ at $x = 10^4$.
\end{proof}

\begin{remark}[A heuristic parabolic multifractal parametrization]
If one \emph{postulates} a Gaussian (log-normal) cascade approximation for the normalized continuous binomial density on dyadic scales, the $L^q$-exponents take the parabolic form $\tau(q) = (q-1)\ln 2 - \frac{\sigma_0^2}{2}q^2$. Its Legendre transform is $f(\alpha) = \ln 2 - \frac{(\alpha - \ln 2)^2}{2\sigma_0^2}$ for all real $\alpha$. This is not a theorem. No multifractal measure has been constructed from the continuous binomial coefficients, the value $\sigma_0^2 = \frac12$ is chosen rather than derived, and with this normalization $\tau(1) = -\sigma_0^2/2 \neq 0$, which departs from the standard convention $\tau(1) = 0$ for probability measures. Here the exponents are in nats per dyadic scale; the value $D_1 = \tau'(1) = \ln 2 - \frac12$ becomes a dimension in the usual sense only after division by $\ln 2$, giving $1 - \frac{1}{2\ln 2}$. The agreement of $D_1 = \ln 2 - \frac12$ with the entropy defect $\lim_{x\to\infty}(x^2\ln 2 - \mathcal{E}(x))/x^2 = \ln 2 - \frac12$ is not independent evidence: both quantities are fixed by the same value $\int_0^1 H = \frac12$ used to choose $\sigma_0^2$.
\end{remark}

\section[Self-Similar Nodal Lattices outside the Simplex]{Self-Similar Nodal Lattices in \texorpdfstring{$\mathbb{R}^2 \setminus \Delta$}{R\textasciicircum 2 \textbackslash{} Delta}}
% ==============================================================================

The binomial coefficient extends off the integers through Euler's reflection formula,
\begin{equation}
\label{eq:reflection}
\binom{x}{y} = -\frac{1}{\pi}\frac{\sin(\pi y)\sin(\pi(x-y))}{\sin(\pi x)} \cdot \frac{\Gamma(y-x)\Gamma(-y)}{\Gamma(-x)},
\end{equation}
valid for $x \notin \mathbb{Z}$ and $y, x-y \notin \mathbb{Z}_{\ge 0}$ (by the removable-singularity convention it extends continuously to those lines). For $x \notin \mathbb{Z}$ the zeros of $\binom{x}{y}$ lie only on the lines $y \in \mathbb{Z}_{<0}$ and $x - y \in \mathbb{Z}_{<0}$. At the other integer values of $y$ and $x - y$ the zeros of the sines cancel against poles of the Gamma factors; for instance $\binom{0.5}{1} = 0.5$, while $\binom{0.5}{-1} = \binom{0.5}{1.5} = 0$. The two families of nodal lines are each parallel, so by themselves they cut the plane into parallelograms. Triangular chambers appear only when the pole lines $x \in \mathbb{Z}_{<0}$ of $\Gamma(x+1)$ are added. The formula plays no role in the dimension count below. At negative integer upper index the coefficients are $\binom{-a}{j} = (-1)^j\binom{a+j-1}{j}$ for $a \ge 1$, $j \ge 0$, so the parity pattern there is an affine image of the one at positive upper index.

\begin{proposition}[Box-counting dimension of Pascal's triangle modulo $2$; classical]
\label{prop:pascal_mod2}
Let $P_k \coloneqq \{2^{-k}(n, j) : 0 \le j \le n < 2^k,\ \binom{n}{j} \equiv 1 \pmod 2\}$. Then $P_k$ converges in the Hausdorff metric to a set $\mathcal{S}$ affinely equivalent to the Sierpi\'nski gasket, and
\begin{equation}
\dim_{\text{box}}(\mathcal{S}) = \dim_H(\mathcal{S}) = \frac{\ln 3}{\ln 2}.
\end{equation}
The same holds for the rescaled parity pattern of $\binom{-a}{j}$, $1 \le a \le 2^k$, $0 \le j$, $a + j \le 2^k$.
\end{proposition}

\begin{proof}
By Lucas' theorem \cite{lucas1878}, $\binom{n}{j}$ is odd if and only if every binary digit of $j$ is at most the corresponding digit of $n$. Hence the rows $0 \le n < 2^{k}$ split into three blocks, rows $n < 2^{k-1}$ and the two copies of that block translated by $(2^{k-1}, 0)$ and $(2^{k-1}, 2^{k-1})$, with the middle block (all entries even) removed. So $P_k = \bigcup_{i=1}^3 \Phi_i(P_{k-1})$ with the three contractions $\Phi_1(z) = z/2$, $\Phi_2(z) = z/2 + (\frac12, 0)$, $\Phi_3(z) = z/2 + (\frac12, \frac12)$. By Hutchinson's theorem the $P_k$ converge in the Hausdorff metric to the attractor of $\{\Phi_1, \Phi_2, \Phi_3\}$. This attractor is the image of the Sierpi\'nski gasket under the linear map sending the equilateral triangle to the triangle with vertices $(0,0), (1,0), (1,1)$. The open set condition holds (take the interior of that triangle), so the box and Hausdorff dimensions are both equal to the similarity dimension $s$ with $3 \cdot 2^{-s} = 1$, i.e.\ $s = \ln 3/\ln 2$ \cite{falconer2014}. Equivalently, $\#P_k = 3^k$ (checked for $k \le 8$). The statement for negative upper index follows from $\binom{-a}{j} = (-1)^j\binom{a+j-1}{j}$ and the affine change of variables $(a, j) \mapsto (a + j - 1, j)$.
\end{proof}

% ==============================================================================
\section{Conclusion}
% ==============================================================================

The chapter has one new proved result: the asymptotic expansion of the continuous row entropy through the Barnes $G$-function, with remainder $-\frac{1}{180x^2} + \mathcal{O}(x^{-4})$ (Proposition~\ref{thm:multifractal}). The renormalization factor $m+3$ (Proposition~\ref{assump:kigami_decimation}), the spectral dimension of the Sierpi\'nski simplex (Kigami--Lapidus) and the fractal limit of Pascal's triangle modulo $2$ (Proposition~\ref{prop:pascal_mod2}) are classical; we have derived the first and recalled the others. Two links remain open: whether Beta-kernel forms on $(K,\mu)$ converge to Kigami's form (Problem~\ref{thm:kigami_convergence}), and a multifractal measure whose $L^q$-spectrum is governed by the continuous binomial coefficients, for which we give only a heuristic parabolic parametrization.

% Shared declarations block, \input by every chapter before its bibliography.
% Keep in sync with the "Declarations" section of master_book_unified_quantum_gravity.tex.
\section*{Declarations}

\noindent\textbf{Funding.} This research was financed in part by the Coordena\c{c}\~ao de Aperfei\c{c}oamento de Pessoal de N\'ivel Superior -- Brasil (CAPES) -- Finance Code 001.

\medskip\noindent\textbf{Competing interests.} The author declares no competing interests.

\medskip\noindent\textbf{Use of generative AI tools.} Large language model assistants (Anthropic Claude; Google Gemini, used through the Antigravity agent environment) were used extensively in preparing this work: drafting and editing text and \LaTeX{} source; proposing, expanding and revising derivations and proofs under the author's direction; writing the Python verification scripts and the \textsc{Lean 4} files; searching and checking the literature and references; and adversarial auditing of proofs, numerical claims and code. These audits found errors, some of them originally introduced with AI assistance. The errors and their corrections are recorded in the repository file \texttt{WORKPLAN.md}. AI tools are not authors. The author chose the research questions and the overall framework, reviewed the material, and is responsible for the content, including any remaining errors.

\medskip\noindent\textbf{Verification status.} Results labelled Theorem, Proposition or Lemma are proved in the text or cited as classical; statements labelled Conjecture are not proved. The numerical scripts compare formulas of the text with independent computations and are checks, not proofs. The \textsc{Lean 4} modules in \texttt{formal\_proofs\_book/Book} are a naming skeleton and do not verify the mathematics of this chapter.

\medskip\noindent\textbf{Code and data availability.} Sources, numerical scripts and the audit log are available at
\begin{center}\url{https://github.com/reinaldomsilvafilho-netizen/quantum-gravity-lean4}.\end{center}
The monograph is archived on Zenodo, \href{https://doi.org/10.5281/zenodo.22290043}{doi:10.5281/zenodo.22290043}, under the CC~BY~4.0 license.


\begin{thebibliography}{99}
\bibitem{silvafilho2026simplex} R.~M. Silva-Filho, \emph{Analytic Continuation of Pascal's Simplex: Continuous Multinomial Integrals, Polytope Boundary Recurrences, and Simplicial Fractional Calculus}, Chapter~3 of \emph{Geometry, Tensors, and Quantum Gravity}, Universidade Federal de Lavras, 2026.
\bibitem{silvafilho2026waves} R.~M. Silva-Filho, \emph{Nonlinear Simplicial Waves and Anomalous Porous Transport Induced by Beta-Kernel Fractional Laplacians}, Chapter~4 of \emph{Geometry, Tensors, and Quantum Gravity}, Universidade Federal de Lavras, 2026.
\bibitem{silvafilho2026inter} R.~M. Silva-Filho, \emph{Inter-Dimensional Simplicial Transforms, Fractional Boundary Traces, and Matrix Beta-Kernels}, Chapter~5 of \emph{Geometry, Tensors, and Quantum Gravity}, Universidade Federal de Lavras, 2026.
\bibitem{kigamilapidus1993} J. Kigami and M. L. Lapidus, \emph{Weyl's problem for the spectral distribution of Laplacians on p.c.f.\ self-similar fractals}, Communications in Mathematical Physics, vol. 158, no. 1, pp. 93--125, 1993, \href{https://doi.org/10.1007/BF02097233}{doi:10.1007/BF02097233}.
\bibitem{kumagaisturm2005} T. Kumagai and K.-T. Sturm, \emph{Construction of diffusion processes on fractals, $d$-sets, and general metric measure spaces}, Journal of Mathematics of Kyoto University, vol. 45, no. 2, pp. 307--327, 2005, \href{https://doi.org/10.1215/kjm/1250281992}{doi:10.1215/kjm/1250281992}.
\bibitem{kigami2001} J. Kigami, \emph{Analysis on Fractals}, Cambridge Tracts in Mathematics, vol. 143, Cambridge University Press, Cambridge, 2001, \href{https://doi.org/10.1017/cbo9780511470943}{doi:10.1017/cbo9780511470943}.
\bibitem{strichartz2006} R. S. Strichartz, \emph{Differential Equations on Fractals: A New Approach}, Princeton University Press, Princeton, 2006, \href{https://doi.org/10.1515/9780691186832}{doi:10.1515/9780691186832}.
\bibitem{barnes1900} E. W. Barnes, \emph{The Theory of the G-Function}, Quarterly Journal of Pure and Applied Mathematics, vol. 31, pp. 264--314, 1900, \href{https://zbmath.org/?q=an:31.0425.01}{JFM 31.0425.01}.
\bibitem{halsey1986} T. C. Halsey, M. H. Jensen, L. P. Kadanoff, I. Procaccia, B. I. Shraiman, \emph{Fractal measures and their singularities: The characterization of strange sets}, Physical Review A, vol. 33, no. 2, pp. 1141--1151, 1986, \href{https://doi.org/10.1103/PhysRevA.33.1141}{doi:10.1103/PhysRevA.33.1141}.
\bibitem{falconer2014} K. Falconer, \emph{Fractal Geometry: Mathematical Foundations and Applications}, 3rd ed., John Wiley \& Sons, Chichester, 2014, ISBN~978-1-119-94239-9.
\bibitem{lucas1878} E. Lucas, \emph{Th\'eorie des Fonctions Num\'eriques Simplement P\'eriodiques}, American Journal of Mathematics, vol. 1, no. 2, pp. 184--240, 1878, \href{https://doi.org/10.2307/2369308}{doi:10.2307/2369308}.
\end{thebibliography}

\end{document}
